\documentclass[12pt]{article}
\usepackage[T2A]{fontenc}
\usepackage[utf8]{inputenc}
\usepackage{amsmath}
\usepackage{amssymb}
\usepackage[english,russian]{babel}
\usepackage[left=2.5cm, right=2cm, top=2cm, bottom=2cm]{geometry}
\setlength{\parindent}{0pt}
\setlength{\parskip}{0pt}

\begin{document}

\noindent\hspace*{-0.5cm}{\Large\bfseries Вариант 3}

\bigskip

Графы \((G(U(t), E(t), Q(t))_{R^*})_i\), \(i = 1, \dots, 3\), формируются по следующим правилам.

\bigskip

1. Построить матрицы \(\|D_a^{R_1}\| = \|d_{aa}^{R_1}\|\), \(\|D_b^{R_2}\| = \|d_{bb}^{R_2}\|\), \(\|D_c^{R_3}\| = \|d_{cc}^{R_3}\|\) таким образом, чтобы
\[
d_{ij}^{R^k} =
\begin{cases}
1, & \text{если } S_i R_k S_j, \, S_i, S_j \in B_1; \, k = 1, \dots, 3; \\
0, & \text{если } S_i R_k S_j, \, S_i, S_j \in B_1, \, \text{не выполняется},
\end{cases}
\tag{1}
\]
где \(a, b, c\) "--- известные константы.

\bigskip

2. Из оставшихся элементов \(\{S\}_{1}^{R_j}\), \(j = 1, \dots, 3\), каждой матрицы построить квадратные матрицы \(\|E_{R_j}\|_j = \|e_{R_j R_j}\|_j\), \(j = 1, \dots, 3\) таким образом, чтобы
\[
e_{iv} =
\begin{cases}
1, & \text{если } S_i R_j S_v, \, j = 1, \dots, 3, \\
0, & \text{в остальных случаях}.
\end{cases}
\tag{2}
\]

\end{document}